% Converted from README.md
\subsection{Primer greške u radu sa pokazivačima}\label{primer-greux161ke-u-radu-sa-pokazivaux10dima}

U slučaju da postoji putanja programa na kojoj može doći do nepravilnog pristupa memorijskim lokacijama, što je čest slučaj pri radu sa pokazivačima, \texttt{KLEE} će sam prijaviti grešku i generisati ulaz na osnovu koga možemo popraviti svoj program. Primer ovakvog programa se može videti u primeru \texttt{03\_err\_pointer}.

Označićemo simbolički ulaz i uključiti potrebno zaglavlje i zatim sačuvati kod u datoteku \texttt{pointer\_error\_sym.c}. Na osnovu nje ćemo napraviti odgovarajući bitkod i pokrenuti alat KLEE. KLEE je pronašao \texttt{8} putanja i generisao \texttt{5} test primera, a pritom je i javio \texttt{memory\ error:\ out\ of\ bound\ pointer}. Test primer koji dovodi do prijavljene greške će biti označen posebnom datotekom sa ekstenzijom \texttt{.ptr.err}. Primetimo da je broj test primera manji od broja istraženih putanja.

\textbf{BITNO:} \texttt{KLEE} neće generisati sve moguće ulaze koji dovode do iste greške. Ukoliko je tako nešto potrebno, prilikom pokretanja dodati opciju \texttt{-{}-emit-all-errors}:

\begin{minted}{bash}
$ klee --emit-all-errors pointer_error_sym.bc
\end{minted}

Ako uvedemo pomenutu opciju prilikom pokretanja alata KLEE, biće generisano \texttt{8} test primera, od čega \texttt{4} označena greškom za memoriju. Čitanjem dobijenih test primera, možemo ispraviti sve putanje koje su dovele do problema sa pokazivačima.

Pokrenuti \texttt{klee-stats} nad direktorijumom koji je generisan u prvom slučaju i sa dodavanjem opcije za emitovanje svih grešaka. Primetiti da se jedino razlikuje vreme potrošeno na solver dok je pokrivenost koda i grana identična.
